\section{Fundamentos teóricos}

\subsection{Numerários e medidas}

\begin{teorema}[Mudança de numerário]\label{teo:numerario}
Sejam $N$ e $M$ dois numerários (processos de preço positivos) e $\mathbb Q^N$ uma medida sob a qual todo preço dividido por $N$ é martingal. A medida definida por
\begin{equation}
\frac{d\mathbb Q^M}{d\mathbb Q^N}\bigg|_{\mathcal F_T}=\frac{M_T/M_0}{N_T/N_0}
\end{equation}
torna martingal todo preço dividido por $M$. Em particular, para qualquer payoff $H$ em $T$,
\begin{equation}
V_0=N_0\,\mathbb E^{\mathbb Q^N}\Big[\frac{H}{N_T}\Big]=M_0\,\mathbb E^{\mathbb Q^M}\Big[\frac{H}{M_T}\Big]
\end{equation}
\citep{gemanetal1995}.
\end{teorema}

\begin{definicao}[Medida forward]
Tomando como numerário o título de cupom zero $P(t,T)$, obtém-se a \emph{medida $T$-forward} $\mathbb Q^T$. Como $P(T,T)=1$,
\begin{equation}
V_0=P(0,T)\,\mathbb E^{\mathbb Q^T}[H],
\end{equation}
e o desconto sai da esperança mesmo com juros estocásticos. Sob $\mathbb Q^T$, a taxa forward $f(t,T)$ e o preço forward de qualquer ativo para entrega em $T$ são martingais.
\end{definicao}

\subsection{O modelo de Hull-White}

\begin{definicao}[Hull-White de um fator]
Sob a medida neutra ao risco,
\begin{equation}\label{eq:hw}
dr_t=\big(\vartheta(t)-ar_t\big)\,dt+\sigma\,dW_t,
\end{equation}
com $a,\sigma>0$ constantes e $\vartheta(t)$ determinística \citep{hullwhite1990}. Com $\vartheta$ constante, é o modelo de Vasicek; com $a=0$, o de Ho-Lee \citep{vasicek1977,holee1986}.
\end{definicao}

\begin{proposicao}[Ajuste à curva]\label{prop:hw-theta}
O modelo reproduz a curva inicial $P^{\mathrm{mkt}}(0,T)$ para todo $T$ se, e somente se,
\begin{equation}
\vartheta(t)=\partial_tf^{\mathrm{mkt}}(0,t)+af^{\mathrm{mkt}}(0,t)+\frac{\sigma^2}{2a}\big(1-e^{-2at}\big).
\end{equation}
Equivalentemente, $r_t=x_t+\alpha(t)$, com $dx_t=-ax_t\,dt+\sigma\,dW_t$, $x_0=0$, e $\alpha(t)=f^{\mathrm{mkt}}(0,t)+\frac{\sigma^2}{2a^2}\big(1-e^{-at}\big)^2$.
\end{proposicao}

\begin{proposicao}[Preço do título]
Em Hull-White, $P(t,T)=A(t,T)\,e^{-B(t,T)r_t}$, com
\begin{equation}
B(t,T)=\frac{1-e^{-a(T-t)}}{a},\qquad \ln A(t,T)=\ln\frac{P^{\mathrm{mkt}}(0,T)}{P^{\mathrm{mkt}}(0,t)}+B(t,T)f^{\mathrm{mkt}}(0,t)-\frac{\sigma^2}{4a}\big(1-e^{-2at}\big)B(t,T)^2.
\end{equation}
\end{proposicao}

\begin{proof}[Esboço]
$\int_t^Tr_s\,ds$ é gaussiana condicionalmente a $\mathcal F_t$, pois $x$ é um processo de Ornstein-Uhlenbeck. Logo $P(t,T)=\exp\big(-m+\tfrac12s^2\big)$, com $m$ e $s^2$ a média e a variância condicionais da integral; ambas são afins em $r_t$, e o coeficiente de $r_t$ é $B(t,T)$. A parte determinística é fixada pela condição de reproduzir $P^{\mathrm{mkt}}(0,\cdot)$.
\end{proof}

\subsection{Opções sobre títulos e caps}

\begin{teorema}[Opção sobre título de cupom zero]
Uma call europeia de vencimento $T$ e strike $X$ sobre o título que vence em $S>T$ vale
\begin{gather}
\mathrm{ZBC}=P(0,S)\,N(h)-X\,P(0,T)\,N(h-\sigma_P),\\
\sigma_P=\sigma\sqrt{\frac{1-e^{-2aT}}{2a}}\,B(T,S),\qquad h=\frac1{\sigma_P}\ln\frac{P(0,S)}{X\,P(0,T)}+\frac{\sigma_P}{2}.
\end{gather}
\end{teorema}

\begin{proof}
Sob $\mathbb Q^T$, $P(T,S)=P(t,S)/P(t,T)$ avaliado em $T$ é o preço forward, um martingal lognormal com variância $\sigma_P^2$. A fórmula é a de Black para esse preço forward \citep{black1976,brigomercurio2006}.
\end{proof}

\begin{proposicao}[Caplet como put sobre título]
Um caplet de taxa simples $L(T,S)$ com strike $K$, período $\delta=S-T$ e nominal $1$, pago em $S$, vale o mesmo que $(1+K\delta)$ puts europeias de vencimento $T$ e strike $1/(1+K\delta)$ sobre o título $P(\cdot,S)$.
\end{proposicao}

\begin{proof}
O payoff $\delta\,(L-K)^+$ em $S$ vale em $T$ $P(T,S)\,\delta(L-K)^+$. Com $1+\delta L=1/P(T,S)$, isso é $(1-(1+K\delta)P(T,S))^+=(1+K\delta)\big(\tfrac1{1+K\delta}-P(T,S)\big)^+$.
\end{proof}

\begin{teorema}[Jamshidian]
Num modelo de um fator em que $P(T,S_i)$ é decrescente em $r_T$, uma opção sobre um título com cupons $c_i$ é uma carteira de opções sobre títulos de cupom zero, com strikes $X_i=P(T,S_i)|_{r_T=r^*}$, onde $r^*$ resolve $\sum_ic_iP(T,S_i)|_{r^*}=X$ \citep{jamshidian1989}.
\end{teorema}
